% 系统架构图 TikZ 版本
% 用于替代 ch5_architecture.png

\begin{figure}[H]
  \centering
  \begin{tikzpicture}[
    box/.style={
      rectangle,
      rounded corners=2pt,
      draw=black,
      fill=#1,
      minimum width=2.0cm,
      minimum height=0.7cm,
      text centered,
      font=\small
    },
    box/.default={white},
    layerbox/.style={
      rectangle,
      rounded corners=2pt,
      draw=black,
      fill=#1,
      minimum width=1.6cm,
      minimum height=0.6cm,
      text centered,
      font=\footnotesize
    },
    layerbox/.default={white},
    group/.style={
      rectangle,
      draw=black!60,
      dashed,
      rounded corners=3pt,
      inner sep=6pt
    },
    arrow/.style={
      -{Stealth[length=5pt, width=3.5pt]},
      thick,
      black
    },
    bidirarrow/.style={
      {Stealth[length=4pt, width=3pt]}-{Stealth[length=4pt, width=3pt]},
      thick,
      black
    }
  ]
  
  % 顶层：Presentation layer
  \node[box={blue!10}] (frontend) at (0, 0) {\textbf{Presentation layer}\\(React frontend)};
  
  % 第二层：Service layer
  \node[box={green!10}] (backend) at (0, -1.5) {\textbf{Service layer}\\(FastAPI)};
  
  % 第三层：Data layer
  \node[box={orange!10}] (data) at (-2.5, -3) {\textbf{Data layer}\\(SQLite)};
  \node[box={orange!10}] (storage) at (2.5, -3) {\textbf{Storage layer}\\(File system)};
  
  % 底层：Model layer
  \node[layerbox={purple!10}] (drenet) at (-2.5, -4.5) {DRENet};
  \node[layerbox={purple!10}] (yolo) at (-0.8, -4.5) {YOLO};
  \node[layerbox={purple!10}] (fcos) at (0.8, -4.5) {FCOS};
  \node[layerbox={purple!10}] (ensemble) at (2.5, -4.5) {Ensemble};
  
  % 垂直箭头
  \draw[bidirarrow] (frontend.south) -- (backend.north);
  \draw[arrow] (backend.south west) -- (data.north east);
  \draw[arrow] (backend.south east) -- (storage.north west);
  
  % 到模型层的箭头（分叉）
  \draw[arrow] (backend.south) -- ++(0, -0.3) -| (drenet.north);
  \draw[arrow] (backend.south) -- ++(0, -0.3) -| (yolo.north);
  \draw[arrow] (backend.south) -- ++(0, -0.3) -| (fcos.north);
  \draw[arrow] (backend.south) -- ++(0, -0.3) -| (ensemble.north);
  
  % 分组标签
  \node[font=\footnotesize, left] at (-3.5, -4.5) {Model layer};
  
  \end{tikzpicture}
  \caption{系统总体架构图}
  \label{fig:ch5_architecture}
\end{figure}
